% Part: many-valued-logic
% Chapter: three-valued-logics
% Section: multiple-designation

\documentclass[../../../include/open-logic-section]{subfiles}

\begin{document}

\olfileid[id]{mvl}{thr}{mul}

\olsection{Menetapkan Bukan Hanya~$\True$}

Sejauh ini, semua logika yang kita lihat mempunyai himpunan nilai kebenaran
yang ditetapkan $V^+=\{\True\}$; yakni, sesuatu dianggap benar jika dan hanya
jika nilai kebenarannya adalah~$\True$. Akan tetapi, sesuatu dapat pula
dianggap benar asalkan tidak bernilai~$\False$. Dalam hal itu, kita memperoleh
suatu logika dengan menetapkan dalam matriksnya, misalnya,
$V^+=\{\True,\Undef\}$.

\begin{defn}
\emph{Logika paradoks}~$\LogLP$ didefinisikan dengan menggunakan matriks
berikut:
\begin{enumerate}
  \item Bahasa proposisional baku~$\Lang L_0$ dengan $\lnot$, $\land$,
  $\lor$, $\lif$.
  \item Himpunan nilai kebenaran~$V=\{\True,\Undef,\False\}$.
  \item $\True$ dan~$\Undef$ ditetapkan, yakni
  $V^+=\{\True,\Undef\}$.
  \item Fungsi kebenarannya sama seperti dalam logika Kleene kuat.
\end{enumerate}
\end{defn}

\begin{defn}
\emph{Logika ketakbermaknaan} Halld\'en~$\LogHal$ didefinisikan dengan
menggunakan matriks berikut:
\begin{enumerate}
  \item Bahasa proposisional baku~$\Lang L_0$ dengan $\lnot$, $\land$,
  $\lor$, $\lif$, dan suatu konektif $1$-tempat~$+$.
  \item Himpunan nilai kebenaran~$V=\{\True,\Undef,\False\}$.
  \item $\True$ dan~$\Undef$ ditetapkan, yakni
  $V^+=\{\True,\Undef\}$.
  \item Fungsi kebenarannya sama seperti dalam logika Kleene lemah, ditambah
  dengan operator ``tidak bermakna'':
  \begin{center}
    \begin{tabular}{c|c}
    $\tf{+}$ & \\
    \hline
    $\True$ & $\False$ \\
    $\Undef$ & $\True$ \\
    $\False$ & $\False$
    \end{tabular}
  \end{center}
\end{enumerate}
\end{defn}

Berbeda dari logika Kleene yang tabel kebenarannya digunakan oleh kedua logika
ini, kedua logika ini \emph{memang} mempunyai tautologi.

\begin{prop}\ollabel{prop:LP-taut-CL} Tautologi dari~$\LogLP$ sama dengan
  tautologi logika proposisional klasik.
\end{prop}

\begin{proof}
  Berdasarkan \olref[syn][sub]{prop:mvl-cl}, jika
  $\Entails[\LogLP]!A$, maka $\Entails[\LogCL]!A$. Untuk menunjukkan arah
  sebaliknya, kita akan membuktikan bahwa jika terdapat
  !!a{valuation}~$\pAssign v\colon\PVar\to\{\False,\True,\Undef\}$
  sedemikian sehingga $\pValue v(!A)[\LogKs]=\False$, maka terdapat
  !!a{valuation}~$\pAssign{v'}\colon\PVar\to\{\False,\True\}$ sedemikian
  sehingga $\pValue{v'}(!A)[\LogCL]=\False$. Hal ini membuktikan hasil bagi
  $\LogLP$, sebab $\LogKs$ dan~$\LogLP$ mempunyai fungsi kebenaran
  karakteristik yang sama, sedangkan $\False$ merupakan satu-satunya nilai
  kebenaran~$\LogLP$ yang tidak ditetapkan (itulah satu-satunya perbedaan
  antara $\LogLP$ dan~$\LogKs$). Jadi, jika
  $\Entails/[\LogLP]!A$, maka untuk suatu !!{valuation}~$\pAssign v$ berlaku
  $\pValue v(!A)[\LogLP]=\pValue v(!A)[\LogKs]=\False$. Berdasarkan klaim
  yang sedang kita buktikan, $\pValue{v'}(!A)[\LogCL]=\False$; yakni,
  $\Entails/[\LogCL]!A$.

  Untuk membuktikan klaim tersebut, mula-mula kita definisikan $\pAssign{v'}$
  dengan
  \[
    \pAssign {v'}(p) =
    \begin{cases}
    \True & \text{jika } \pAssign {v}(p) \in \{\True, \Undef\}\\
    \False & \text{dalam kasus lainnya.}
  \end{cases}
  \]
  Sekarang kita buktikan dengan induksi pada~$!A$ bahwa (a)~jika
  $\pValue v(!A)[\LogKs]=\False$, maka
  $\pValue{v'}(!A)[\LogCL]=\False$, dan (b)~jika
  $\pValue v(!A)[\LogKs]=\True$, maka
  $\pValue{v'}(!A)[\LogCL]=\True$.
  \begin{enumerate}
    \item Dasar induksi: $!A\ident p$. Berdasarkan
    \olref[syn][val]{defn:pValue}, jika $\pValue v(!A)[\LogKs]=\False$,
    maka $\pAssign v(p)=\False$, sehingga menurut definisi~$\pAssign{v'}$,
    $\pValue{v'}(!A)[\LogCL]=\pAssign{v'}(p)=\False$. Jika
    $\pValue v(!A)[\LogKs]=\True$, maka $\pAssign v(p)=\True$, sehingga
    $\pValue{v'}(!A)[\LogCL]=\pAssign{v'}(p)=\True$. Jadi, (a) dan~(b)
    keduanya berlaku.

    Untuk langkah induksi, bahas kasus-kasus berikut:
    \item $!A\ident\lnot !B$.
    \begin{enumerate}
      \item Andaikan $\pValue v(\lnot!B)[\LogKs]=\False$. Menurut definisi
      $\tf{\lnot}[\LogKs]$, $\pValue v(!B)[\LogKs]=\True$. Berdasarkan
      hipotesis induksi dalam kasus~(b), kita memperoleh
      $\pValue{v'}(!B)[\LogCL]=\True$, sehingga
      $\pValue{v'}(\lnot !B)[\LogCL]=\False$.
      \item Andaikan $\pValue v(\lnot!B)[\LogKs]=\True$. Menurut definisi
      $\tf{\lnot}[\LogKs]$, $\pValue v(!B)[\LogKs]=\False$. Berdasarkan
      hipotesis induksi dalam kasus~(a), kita memperoleh
      $\pValue{v'}(!B)[\LogCL]=\False$, sehingga
      $\pValue{v'}(\lnot !B)[\LogCL]=\True$.
    \end{enumerate}

    \item $!A\ident(!B\land !C)$.
    \begin{enumerate}
      \item Andaikan $\pValue v(!B\land !C)[\LogKs]=\False$. Menurut
      definisi $\tf{\land}[\LogKs]$, $\pValue v(!B)[\LogKs]=\False$ atau
      $\pValue v(!C)[\LogKs]=\False$. Berdasarkan hipotesis induksi dalam
      kasus~(a), kita memperoleh $\pValue{v'}(!B)[\LogCL]=\False$ atau
      $\pValue{v'}(!C)[\LogCL]=\False$, sehingga
      $\pValue{v'}(!B\land !C)[\LogCL]=\False$.
      \item Andaikan $\pValue v(!B\land !C)[\LogKs]=\True$. Menurut definisi
      $\tf{\land}[\LogKs]$, $\pValue v(!B)[\LogKs]=\True$ dan
      $\pValue v(!C)[\LogKs]=\True$. Berdasarkan hipotesis induksi dalam
      kasus~(b), kita memperoleh $\pValue{v'}(!B)[\LogCL]=\True$ dan
      $\pValue{v'}(!C)[\LogCL]=\True$, sehingga
      $\pValue{v'}(!B\land !C)[\LogCL]=\True$.
    \end{enumerate}
  \end{enumerate}

  Kedua kasus lainnya serupa dan ditinggalkan sebagai latihan. Sebagai
  alternatif, pembuktian di atas menetapkan hasil bagi semua !!{formula}s yang
  hanya memuat $\lnot$ dan~$\land$. Selanjutnya kita dapat menggunakan fakta
  bahwa, baik dalam $\LogKs$ maupun~$\LogCL$, bagi setiap $\pAssign v$ berlaku
  $\pValue v(!B\lor !C)=\pValue v(\lnot(\lnot!B\land\lnot !C))$ dan
  $\pValue v(!B\lif !C)=\pValue v(\lnot(!B\land\lnot !C))$.
\end{proof}

\begin{prob}
  Lengkapi pembuktian \olref[mvl][thr][mul]{prop:LP-taut-CL}; yakni,
  buktikan (a) dan~(b) untuk kasus $!A\ident(!B\lor !C)$ dan
  $!A\ident(!B\lif !C)$.
\end{prob}

\begin{prob}
Buktikan bahwa setiap tautologi klasik merupakan tautologi dalam~$\LogHal$.
\end{prob}

Meskipun mempunyai tautologi yang sama dengan logika klasik, kedua logika ini
mempunyai relasi konsekuensi yang berbeda. Sebagai contoh, $\LogLP$ bersifat
\emph{parakonsisten}, sebab $\lnot p,p\Entails/ q$; jadi, prinsip ledakan
$\lnot !A,!A\Entails !B$ tidak berlaku secara umum. (Prinsip itu berlaku pada
beberapa kasus $!A$ dan~$!B$, misalnya jika $!B$ merupakan tautologi.)

\begin{prob}
  Manakah di antara relasi-relasi berikut yang berlaku dalam (a)~$\LogLP$ dan
  (b)~$\LogHal$? Berikan tabel kebenaran bagi masing-masing relasi.
  \begin{enumerate}
    \item $p,p\lif q\Entails q$
    \item $\lnot q,p\lif q\Entails\lnot p$
    \item $p\lor q,\lnot p\Entails q$
    \item $\lnot p,p\Entails q$
    \item $p\Entails p\lor q$
    \item $p\lif q,q\lif r\Entails p\lif r$
  \end{enumerate}
\end{prob}

Apa yang terjadi jika $\Undef$ ditetapkan dalam~$\LogLuk[3]$?

\begin{defn}
  \emph{R-Mingle bernilai~$3$}~$\LogRM[3]$ didefinisikan dengan menggunakan
  matriks berikut:
  \begin{enumerate}
    \item Bahasa proposisional baku~$\Lang L_0$ dengan $\lfalse$, $\lnot$,
    $\land$, $\lor$, $\lif$.
    \item Himpunan nilai kebenaran~$V=\{\True,\Undef,\False\}$.
    \item $\True$ dan~$\Undef$ ditetapkan, yakni
    $V^+=\{\True,\Undef\}$.
    \item Fungsi kebenarannya sama seperti dalam logika
    \L ukasiewicz~$\LogLuk[3]$.
  \end{enumerate}
\end{defn}

\begin{prob}
  Manakah di antara relasi-relasi berikut yang berlaku dalam~$\LogRM[3]$?
  \begin{enumerate}
    \item $p,p\lif q\Entails q$
    \item $p\lor q,\lnot p\Entails q$
    \item $\lnot p,p\Entails q$
    \item $p\Entails p\lor q$
  \end{enumerate}
\end{prob}

Tabel kebenaran yang berbeda kadang-kadang dapat menghasilkan logika yang sama
(relasi perikutan yang sama) hanya dengan mengubah nilai-nilai yang ditetapkan.
Sebagai contoh, hal ini terjadi jika dalam logika G\"odel kita menggunakan
$V^+=\{\True,\Undef\}$ alih-alih~$\{\True\}$.

\begin{prop}\ollabel{prop:gl-udes}
  Matriks dengan $V=\{\False,\Undef,\True\}$,
  $V^+=\{\True,\Undef\}$, dan fungsi kebenaran logika G\"odel bernilai~$3$
  mendefinisikan logika klasik.
\end{prop}

\begin{proof}
  Latihan.
\end{proof}

\begin{prob}
  Buktikan \olref[mvl][thr][mul]{prop:gl-udes} dengan menunjukkan bahwa bagi
  logika~$\Log L$ yang didefinisikan sama seperti logika G\"odel tetapi dengan
  $V^+=\{\True,\Undef\}$, jika $\Gamma\Entails/[\Log L]!B$, maka
  $\Gamma\Entails/[\LogCL]!B$. Gunakan gagasan dalam
  \olref[mvl][thr][mul]{prop:LP-taut-CL}; tetapi alih-alih membuktikan sifat
  (a) dan~(b), tunjukkan bahwa $\pValue v(!A)[\LogGod]=\False$ jika dan hanya
  jika $\pValue{v'}(!A)[\LogCL]=\False$ (dan karena itu
  $\pValue v(!A)[\LogGod]\in\{\True,\Undef\}$ jika dan hanya jika
  $\pValue{v'}(!A)[\LogCL]=\True$). Jelaskan mengapa hal ini membuktikan
  proposisi tersebut.
\end{prob}

\end{document}
